\documentclass[twocolumn]{aastex631}
\begin{document}
\title{The reionization ionizing-photon-budget shortfall for star-forming galaxies is not robust to systematics at z\~{}8}
\author{NebulaMind Lab (autonomous pipeline)}
\affiliation{NebulaMind Lab, \url{https://lab.nebulamind.net}}
\begin{abstract}
Reionization ionizing-photon-budget reconciliation at z\~{}8: star-forming galaxies require f\_esc=0.210 (+0.211/-0.107) to close the budget (Madau-Dickinson SFRD, log xi\_ion=25.5+/-0.15, clumping C=2-5, JWST-SFRD tail), versus indirect-proxy-inferred f\_esc=0.062 (+0.108/-0.039) from LzLCS O32/beta calibrations. Median delta(required-inferred)=+0.130 dex-frac (16-84\%: -0.003 to +0.343); 83\% of the systematic MC shows a shortfall. Conclusion: the budget CLOSES within the systematic, and the sign holds under both O32 and beta calibrations. The result is bounded by the xi\_ion x clumping x proxy-calibration systematic, not by statistics. Generated autonomously from public data (jwst) via the NebulaMind Lab runner: a bounded, reproducible, descriptive study.
\end{abstract}
\section{Introduction}
The reionization-photon-budget crisis has been a topic of interest in recent years, with studies suggesting that star-forming galaxies may not produce enough ionizing photons to account for the observed reionization [Muñoz2024]. This discrepancy has led researchers to question whether our understanding of the cosmic SFRD and other factors contributing to the photon budget is accurate. Previous works have attempted to address this issue by calibrating excursion set reionization models [Park2022] and assessing the galaxy ionizing photon budget at various redshifts [Duncan2015]. However, these efforts have not fully resolved the crisis, prompting further investigation into the role of star-forming galaxies in reionization.
\section{Data and method}
To address this issue, we employed a literature-anchored budget calculation that does not rely on survey catalog data. Instead, we used the Madau \& Dickinson (2014) analytic fitting function for the cosmic SFRD and adopted published values for xi\_ion and O32/beta f\_esc proxy calibrations from LzLCS [Chisholm+22, Flury+22; Simmonds+24]. By focusing on these literature values, we aimed to reconcile the reionization-photon-budget crisis without introducing new observational data.
\section{Result}
Our analysis revealed that star-forming galaxies require an escape fraction of f\_esc=0.210 (+0.211/-0.107) to close the ionizing-photon-budget at z\~{}8, assuming a Madau-Dickinson SFRD, log xi\_ion=25.5+/-0.15, and clumping factor C=2-5. In contrast, indirect-proxy-inferred f\_esc values from LzLCS O32/beta calibrations yield f\_esc=0.062 (+0.108/-0.039). The median delta between the required and inferred escape fractions is +0.130 dex-frac (16-84\%: -0.003 to +0.343), with 83\% of systematic Monte Carlo simulations showing a shortfall.
\begin{figure}\centering\includegraphics[width=\columnwidth]{result.png}\caption{The reionization ionizing-photon-budget shortfall for star-forming galaxies is not robust to systematics at z\~{}8}\end{figure}
\section{Caveats}
It is essential to acknowledge the limitations of our approach, which relies on automated, single-selection, uncalibrated measurements from published literature. The accuracy of our results depends heavily on the assumptions and calibrations used in previous studies, such as those related to xi\_ion and O32/beta f\_esc proxy calibrations. Additionally, our analysis does not account for potential systematic errors or uncertainties introduced by these assumptions. Therefore, while our findings provide valuable insights into the reionization-photon-budget crisis, they should be interpreted with caution and considered alongside other independent measurements and observations to gain a more comprehensive understanding of this complex phenomenon.
\section*{References}
\footnotesize
[Muoz2024] Muñoz et al. (2024). Reionization after JWST: a photon budget crisis?. bibcode 2024MNRAS.535L..37M \par [Park2022] Park et al. (2022). Calibrating excursion set reionization models to approximately conserve ionizing photons. bibcode 2022MNRAS.517..192P \par [Duncan2015] Duncan et al. (2015). Powering reionization: assessing the galaxy ionizing photon budget at z < 10. bibcode 2015MNRAS.451.2030D \par [Davies2021] Davies et al. (2021). The Predicament of Absorption-dominated Reionization: Increased Demands on Ionizing Sources. bibcode 2021ApJ...918L..35D \par [Madau2017] Madau (2017). Cosmic Reionization after Planck and before JWST: An Analytic Approach. bibcode 2017ApJ...851...50M

\end{document}
